\documentclass[12pt]{article}

\newif\ifsol
\soltrue

% Unicode compatibility
\usepackage{iftex}
\ifPDFTeX
  \usepackage[utf8]{inputenc}
  \usepackage[noTeX]{mmap}
  \usepackage[T1]{fontenc}
\fi
% XeTeX does not support mmap
\ifLuaTeX
  \usepackage{luatex85}
  \usepackage[noTeX]{mmap}
\fi

\setlength{\textheight}{9in}
\setlength{\textwidth}{7.1in}
\setlength{\evensidemargin}{-0.2in}
\setlength{\oddsidemargin}{-0.2in}
\setlength{\headsep}{30pt}
\setlength{\topmargin}{-0.3in}

\usepackage{amsthm,amsfonts,amsmath,xspace,tikz,varwidth,cancel,soul}
\usepackage{hyperref,verbatim}
\newtheorem{theorem}{Theorem}
\theoremstyle{definition}
\newtheorem{definition}{Definition}

\newcommand{\addmedskip}{\addvspace{\medskipamount}}

\newcommand{\addbigskip}{\addvspace{\bigskipamount}}

\pagestyle{myheadings}
\markboth{BU CAS CS 538.}{BU CAS CS 538.} 
\newcounter{problemnum}
\setcounter{problemnum}{0}
\newenvironment{problem}
     {\addbigskip \setcounter{partnum}{0}
      \noindent\stepcounter{problemnum}\textbf{Problem
                                             \arabic{problemnum}.\ }}
     {\par\addbigskip}
     
\ifsol
\newenvironment{solution}
     {\addbigskip
      \noindent\color{blue}\textbf{Solution.\ }}
     {\par\addbigskip}
\else
\newenvironment{solution}
     {\comment}
     {\endcomment}
\fi

\newcounter{partnum}
\setcounter{partnum}{0}
\newenvironment{ppart}
     {\addmedskip
      \noindent\stepcounter{partnum}\textbf{(\alph{partnum})}\ }
     {\par\addbigskip}

\newcommand{\Enc}{\mathrm{Enc}}
\newcommand{\Dec}{\mathrm{Dec}}
\newcommand{\eqdef}{\buildrel{\scriptstyle{\mathit{def}}}\over{=}}
\newcommand{\rfrom}{\buildrel{\scriptstyle{\mathrm{R}}}\over{\leftarrow}}

\newcommand{\mbmod}{\,\%\,}

\newcommand{\A}{\mathcal{A}}
\renewcommand{\L}{\mathcal{L}}
\newcommand{\Z}{\mathbb{Z}}
\newcommand{\G}{\mathbb{G}}
\newcommand{\X}{\mathcal{X}}
\newcommand{\Y}{\mathcal{Y}}
\newcommand{\K}{\mathcal{K}}
\newcommand{\M}{\mathcal{M}}
\newcommand{\E}{\mathcal{E}}
\newcommand{\C}{\mathcal{C}}
\newcommand{\R}{\mathcal{R}}
\renewcommand{\S}{\mathcal{S}}
\newcommand{\I}{\mathcal{I}}
\newcommand{\T}{\mathcal{T}}
\newcommand{\randk}{\mathsf k}
\newcommand{\randm}{\mathsf m}
\newcommand{\randc}{\mathsf c}
\newcommand{\Adv}{\mathcal{A}}
\newcommand{\B}{\mathcal{B}}
\newcommand{\SSadv}{\textrm{SS}\textsf{adv}}
\newcommand{\PRGadv}{\textrm{PRG}\textsf{adv}}
\newcommand{\threebitadv}{\textrm{3Bit}\textsf{adv}}
\newcommand{\PRFadv}{\textrm{PRF}\textsf{adv}}
\newcommand{\MACadv}{\textrm{MAC}\textsf{adv}}
\newcommand{\CPAadv}{\textrm{CPA}\textsf{adv}}
\newcommand{\BCadv}{\textrm{BC}\textsf{adv}}
\newcommand{\CRadv}{\textrm{CR}\textsf{adv}}
\newcommand{\DDHadv}{\textrm{DDH}\textsf{adv}}
\newcommand{\comadv}{\textrm{Com}\textsf{adv}}
\newcommand{\pk}{\mathit{pk}}
\newcommand{\sk}{\mathit{sk}}
\newcommand{\qrn}{\mathit{QR}_n}
\newcommand{\qrp}{\mathit{QR}_p}
\newcommand{\qrpo}{\mathit{QR}_{p_1}}
\newcommand{\qrpt}{\mathit{QR}_{p_2}}
\newcommand{\crt}{\mathrm{crt}}




\newcommand{\indist}{  \mathrel{\vcenter{\offinterlineskip
  \hbox{$\sim$}\vskip-.35ex\hbox{$\sim$}\vskip-.35ex\hbox{$\sim$}}}}

\newcommand{\samples}{\xleftarrow{R}}

\newcommand{\hdl}{H_{\mathrm{dl}}}


\begin{document}
\begin{center}
\Large{\textbf{CAS CS 538.  \ifsol Solutions to \fi Problem Set 10}}\\
\smallskip
\large{\textbf{Due electronically via Gradescope, Wednesday April 22, 2026 11:59pm}}\\
\large{\textbf{Om Khadka, U51801771}}
\end{center}

\noindent
\textbf{How to submit:}  See instructions on Problem Set 1. 
 
\section{Commitments}


\begin{problem}[El Gamal commitments]
    We now consider a different commitment than the one from Discussion 10. 
    Let $\G$ again be a group of prime order $q$ generated by $g$. 
    Consider the following protocol between Alice and Bob. 
    Given a message $m\in \G$, Alice runs the following algorithm:
    \begin{itemize}
        \item \underline{\smash{Commit$(m)$}}: pick random $u\samples \G$,  $r\samples \Z_q$ and let $c=(u, g^r, u^r \cdot m)$.
    \end{itemize}
    Alice sends $c$ to Bob and keeps $r$ to herself.
    Bob makes sure that he got three elements of $\G$.
    
    Some time later, Bob wants to know what $m$ is and Alice is ready to reveal it.
    Alice sends $r$ and $m$ to Bob.
    Of course, she may not be honest, so we will call what Bob receives $m'$ and $r'$.
    Bob runs the following algorithm:
    \begin{itemize}
        \item \underline{\smash{Verify$(c, m',r')$}}: 
        First, make sure that $m' \in \G$ and $r' \in \Z_q$. 
        Let $c=(u, v, z)$. Check that $v=g^{r'}$ and that $z=u^{r'} m'$. 
        If so, accept $m'$ as the correct message. Else, reject.
    \end{itemize}
    Note that this commitment scheme is just the encryption scheme known as 
    multiplicative ElGamal (Exercise 11.5), but without decryption.\footnote{
    In fact, a similar commitment scheme can be built out of any encryption 
    scheme that uses randomness $r$: to commit to $m$, output $\pk$ and $E(\pk, m)$ 
    where $E$ uses randomness $r$ and no other randomness. To verify a revealed 
    message $m'$, use $r'$ to re-encrypt $m'$ under $\pk$ and check that the result 
    matches the received commitment. It works as long as for every public key --- 
    even adversarially generated --- correct decryption is unique.}

    \begin{ppart}
        Prove that this commitment scheme is \emph{computationally hiding} 
        under the DDH assumption.
        
        The definition of computationally hiding commitments is essentially 
        the same as the definition of semantically security encryption schemes: 
        given $\G, g$, the adversary outputs $m_0$ and $m_1$, 
        gets a commitment to $m_b$ for random bit $b$ and outputs $\hat{b}$, 
        with the usual measure of adversarial advantage 
        $\comadv = \lvert\Pr[\hat{b}=1\, |\,b=0]-\Pr[\hat{b}=1\,|\,b=1]\rvert$.
    \end{ppart}

\begin{solution}
This will be proven by reducing the game down to DDH.\\

\noindent Assume, for contradiction, that some eff. adv. $\mathcal{A}$ exists, that can crack the hiding property with non-negligible advantage $\epsilon$.
It's important to note that, from the textbook (Boneh-Shoup: 2.2.5), $\mathcal{A}$'s advantage in the bit-guessing version of this game would then be $\epsilon^* = \epsilon/2$, where $\epsilon^* = \lvert \Pr[\hat{b}=b] - 1/2 \rvert$.\\

\noindent Now suppose another adv, $\mathcal{B}$. $\mathcal{B}$ will use $\mathcal{A}$ to break DDH. $\mathcal{B}$ will be attempting to directly crack challenger $\mathcal{C}_{\text{DDH}}$:
\begin{itemize}
    \item $\mathcal{B}$ will get $(u, v, w) \in \mathbb{G}^3$ from $\mathcal{C}_{\text{DDH}}$.
    \begin{itemize}
      \item In Exp 0, $(u,v,w) = (g^\alpha, g^\beta, g^{\alpha\beta})$.
      \item In Exp 1, $(u,v,w) = (g^\alpha, g^\beta, g^\gamma)$ for independent and rnd $\alpha,\beta,\gamma \in \mathbb{Z}_q$.
    \end{itemize}
    \item $\mathcal{B}$ will run $\mathcal{A}$ on input $(\mathbb{G}, g)$. $\mathcal{A}$ should output $m_0, m_1 \in \mathbb{G}$.
    \item $\mathcal{B}$ will then choose a rnd bit $b \xleftarrow{R} \{0,1\}$.
    \item $\mathcal{B}$ will then make the candidate commitment $c = (u, v, w \cdot m_b)$ and then send it over to $\mathcal{A}$.
    \begin{itemize}
      \item Note: Since $\mathbb{G}$ is a prime-order cyclic group generated by $g$, getting $u \xleftarrow{R} \mathbb{G}$ is (statistically) the same thing as getting $g^\alpha$ for a rnd $\alpha$. This matches the commitment's distribution.
    \end{itemize}
    \item $\mathcal{A}$ returns $\hat{b}$. If $\hat{b} = b$, $\mathcal{B}$ outputs $1$. Otherwise it'll output $0$. \newline
\end{itemize}

\noindent\textbf{Analysis}\\
\noindent Now note, the main difference in the two experiments is in what $\mathcal{C}_{\text{DDH}}$ runs:
\begin{itemize}
    \item[\textbf{Exp 0}] $\mathcal{C}_{\text{DDH}}$ will run $w = g^{\alpha\beta}$. 
    Then, $c = (g^\alpha, g^\beta, (g^\alpha)^\beta \cdot m_b) = (u, g^r, u^r \cdot m_b)$ with $r=\beta$. 
    This is just a valid commitment $m_b$. 
    So, $\mathcal{A}$ is playing with the real hiding game.
    And by definition of bit-guessing advantage:
    \begin{itemize}
      \item $\Pr[\mathcal{B} \to 1 \mid \text{Exp } 0] = \Pr[\hat{b}=b] = \frac{1}{2} + \epsilon^*$
    \end{itemize}
    \item[\textbf{Exp 1}] Now, $\mathcal{C}_{\text{DDH}}$ is running $w = g^\gamma$ with it being uniform in $\mathbb{G}$.
    Note that, in this case, the third component $w \cdot m_b$ will also be uniformly distributed over $\mathbb{G}$ (since multiplication by a fixed group element is a bijection on $\mathbb{G}$).
    Consequently, $c$ is independant of $b$.
    This means that $\mathcal{A}$ can do no better than just rnd guessing:
    \begin{itemize}
      \item $\Pr[\mathcal{B} \to 1 \mid \text{Exp } 1] = 1/2$ \newline
    \end{itemize}
  \end{itemize}

\noindent  $\mathcal{B}$'s DDH advantage can then be calculcated as such:
\[
\text{DDHAdv}[\mathcal{B}, \mathbb{G}] = \lvert \Pr[\mathcal{B} \to 1 \mid \text{Exp } 0] - \Pr[\mathcal{B} \to 1 \mid \text{Exp } 1] \rvert = \left\lvert \left(\frac{1}{2} + \epsilon^*\right) - \frac{1}{2} \right\rvert = \epsilon^* = \frac{\epsilon}{2}.
\]
Since $\epsilon$ is non-negligible, $\text{DDHAdv}[\mathcal{B}, \mathbb{G}]$ is also non-negligible.
This contradicts the DDH assumption, and therefore $\mathcal{A}$ can't exist.
This proves how this commitment scheme is computationally hiding.
\end{solution}

\begin{ppart}
Prove that the commitment is \emph{perfectly binding}: namely, no matter how clever or computationally unbounded Alice is, for every $c$ and every $r_1, r_2$, there is at most one $m$ that Bob will accept.
  In other words, for every $c, r_1, r_2, m_1, m_2$, if Verify$(c, m_1, r_1)$ accepts and Verify$(c, m_2, r_2)$ also accepts, then $m_1 = m_2$.
\end{ppart}

\begin{solution}
This can be proven just by utilizing some algebraic properties (especially of prime-order groups).\\

\noindent Suppose an adv. makes a commitment $c = (u, v, z)$ and claims two valid openings $(m_1, r_1)$ and $(m_2, r_2)$.
For both to be accepted by the Verify algorithm, they have to satisfy the following verification equations:
\begin{itemize}
\item For $(m_1, r_1)$: $v = g^{r_1}$ and $z = u^{r_1} m_1$.
\item For $(m_2, r_2)$: $v = g^{r_2}$ and $z = u^{r_2} m_2$.\\
\end{itemize}

\noindent\textbf{Uniqueness of $r$:}
Note the first components, $g^{r_1} = v = g^{r_2}$.
Since $\G$ is a cyclic group of prime order $q$ generated by $g$, the mapping of $r \mapsto g^r$ is a bijection from $\Z_q$ to $\G$.
Then,\\
$g^{r_1} = g^{r_2} \implies r_1 \equiv r_2 \pmod q$.
And because both $r_1, r_2 \in \Z_q$, we must then have $r_1 = r_2$.\\

\noindent\textbf{Uniqueness of $m$:}
Now we can sub $r_1 = r_2$ into the second components:
\[ z = u^{r_1} m_1 \quad \text{and} \quad z = u^{r_2} m_2 = u^{r_1} m_2. \]
Equating the two expressions for $z$ gives $u^{r_1} m_1 = u^{r_1} m_2$.
Since $\G$ is a group, every element has a unique inverse.
Multiplying both sides by $(u^{r_1})^{-1}$ will cancel the term, resulting in $m_1 = m_2$.\\

\noindent This proves that any commitment, $c$, will have AT MOST one message, $m$, that will pass verification. 
This holds unconditionally, proving that this scheme is perfectly binding.
\end{solution}
\end{problem}

\begin{problem} (20 points)
The above scheme is \emph{perfectly} binding and \emph{computationally} hiding. In discussion 10 we covered Pedersen commitments, which are \emph{computationally} biding and \emph{perfectly} hiding. 

To recap the discussion, a commitment scheme is \emph{perfectly hiding} when
for every $g, h, m_0$ and $m_1$, the distribution of outputs 
produced by Commit$(m_0)$ is identical to the distribution of outputs 
produced by Commit$(m_1)$. (Note that this protects Alice during the commit phase because Bob learns nothing about $m$, regardless of Bob's computational power or any computational assumptions.)

Prove that it's not possible to have a commitment scheme to be both perfectly binding and perfectly hiding.
\end{problem}

\begin{solution}
To prove this question, one could just make a counting argument.
This will touch into Shannon's impossibility theorem for perfect secrecy (and more broadly, the Pigeonhole Theorem).\\

\noindent Suppose, for the sake of contradiction, that there exists a commitment scheme with msg space $\M$, rnd space $\mathcal{R}$, 
and commit space $\C$ that is BOTH perfectly hiding and perfectly binding.
We also assume that $|\M| \ge 2$ (it'd be trivial otherwise).
The Commit algorithm picks $r \xleftarrow{R} \mathcal{R}$ uniformly at rnd.\\

\noindent For any msg $m \in \M$ and commit $c \in \C$, let $N_m(c)$ denote the num. of randomness values $r \in \mathcal{R}$ such that $\text{Commit}(m; r) = c$.
Now, since $r$ is uniform, the probability of producing $c$ from $m$ is $\Pr[\text{Commit}(m) = c] = \frac{N_m(c)}{|\mathcal{R}|}$.\\

\noindent\textbf{Consequence of Perfect Hiding:}
Perfect hiding requires that $\forall m_0, m_1 \in \M$ and $\forall c \in \C$, the distributions are the same:
\[
\Pr[\text{Commit}(m_0) = c] = \Pr[\text{Commit}(m_1) = c] \implies \frac{N_{m_0}(c)}{|\mathcal{R}|} = \frac{N_{m_1}(c)}{|\mathcal{R}|} \implies N_{m_0}(c) = N_{m_1}(c).
\]
Therefore, $N_m(c)$ is independent of $m$.
Now, let $N(c)$ denote this common value.
If $N(c) > 0$, then for \emph{every} $m \in \M$, there exists AT LEAST one randomness $r$ that opens $c$ to $m$.\\

\noindent\textbf{Consequence of Perfect Binding:}
Now perfect binding requires that for any commitment $c$, there's at most one message, $m$, that Bob will accept.
Basically, if $\text{Verify}(c, m_1, r_1)$ and $\text{Verify}(c, m_2, r_2)$ both accept, then it should be the case that$m_1 = m_2$.
This then implies that for any fixed $c$, there can be valid openings for AT LEAST one distinct message, $m$.\\

\noindent\textbf{End:} Now, consider any $c \in \C$ that is in the support of the commitment distribution (this $c$ has to exist, as commitments are generated).
For this $c$, we will have $N(c) > 0$:
\begin{itemize}
  \item By the Consequence of Perfect Hiding, $N(c) > 0$ implies that $c$ can be validly opened to EVERY $m \in \M$. 
  \item By the Consequence of Perfect Binding, $c$ can be validly opened to AT MOST one $m \in \M$.
\end{itemize}
NOw, these two conditions can only work if and only if $|\M| \le 1$.
However, this contradicts the initial assumption that the scheme supports AT LEAST two distinct messages. \\

\noindent Therefore, no commitment scheme can be both perfectly hiding and perfectly binding when $|\M| \ge 2$.
\end{solution}



\section{Weak PRF to PRF}

Recall that in the attack game for a PRF $F$ defined over $(\K, \X, \Y)$,
the adversary $\Adv$ submits a sequence of queries to the challenger, 
and the challenger picks a random key $k$ and answers each query $x_i$ with $F(k, x_i)$ in Experiment 0,
and it picks a random function $f$ from the set of all functions from $\X$ to $\Y$ and answers each query 
$x_i$ with $f(x_i)$ in Experiment 1.
Note that the adversary has full control over the queries.

We can, however, define a weaker notion of a PRF, called a \emph{weak PRF}, 
where the adversary is only allowed to make function queries on random inputs.
More specifically, in the weak PRF game, whenever the adversary $\Adv$ makes a query,
the \emph{challenger} picks a random $x_i \in \X$ and answers the query with $(x_i, F(k, x_i))$ in Experiment 0, 
and with $(x_i, f(x_i))$ in Experiment 1.
As expected, the adversary's advantage in this game, denoted $\mathrm{wPRF}\textsf{adv}[\Adv, F]$, 
is defined as the difference between the probability that $\Adv$ outputs 1 in Experiment 0 and the probability that $\Adv$ outputs 1 in Experiment 1; 
and the PRF $F$ is said to be a weak PRF if $\mathrm{wPRF}\textsf{adv}[\Adv, F]$ is negligible for all 
efficient adversaries $\Adv$. (See definition 4.3 in Section 4.4.1 of the textbook if you want more explanations.)

\begin{problem} (35 points)
Problem 11.1 in the textbook. Correction to the problem: let $q$ be the total number of queries that $\Adv$ makes. Then
$$\mathrm{PRF}^{\mathrm{ro}}\mathsf{adv}[\mathcal{A}, \mathcal{F} ] \le \mathrm{wPRF}\mathsf{adv}[\mathcal{B}, \mathcal{F'}]$$
should be instead
$$\mathrm{PRF}^{\mathrm{ro}}\mathsf{adv}[\mathcal{A}, \mathcal{F} ] \le \mathrm{wPRF}\mathsf{adv}[\mathcal{B}, \mathcal{F'}]+\frac{q^2}{2\cdot |\mathcal{X}|}$$

\paragraph{Interpretation:} We know that $\Adv$ is polytime bounded and hence it can make only a polynomial number of queries. Thus, $q^2$ is negligible if $|\mathcal{X}|$ is exponentially large. So this weak-PRF-to-PRF conversion works for large $|\mathcal{X}|$ (something like $256$-bit is reasonable) but not for small $\mathcal{X}$. 


\paragraph{Approach:}
	Let $\Adv$ be a PRF adversary against $F$. Because we are in the random oracle model, $\Adv$ gets to ask two types of queries: 
\begin{itemize}
	\item $q_p$ PRF queries $m_i$ (for $1\le i \le q_p$) to the PRF (in Game 0) or random function (in Game 1)
	\item $q_H$ random oracle (RO) queries $\hat{m}_i$ (for $1\le i\le q_H$) to $H$	
\end{itemize}
The total number of queries is $q=q_p+q_H$.
	
	Let $p_0$ be the probability that $\Adv$ outputs $1$ in experiment 0 of the PRF game (queries answered by $F$), and $p_1$ be the probability that $\Adv$ outputs $1$ in experiment 1 of the PRG game(queries answered by a random $f$). 
	
	Let $\C$ be the following challenger: pick a random function $H: \M \to\X$ and a random function $f':\X\to \Y$;	answer PRF queries $m_i$ with $f'(H(m_i))$ and RO queries $\hat{m}_i$ with $H(\hat{m}_i)$. Let $p_\C$ be the probability that $\Adv$ outputs $1$ when interacting with $\C$. 
	
	In the text two parts, you will prove $|p_0-p_c|$ and $|p_c-p_1|$ are both small, giving you that $|p_0-p_1|$ is small.

	\begin{ppart}
	Prove that $|p_0-p_c|\le \mathrm{wPRF}\mathsf{adv}[\mathcal{B}, \mathcal{F'}]$ by building an adversary $\B$ that uses $\Adv$ and plays the wPRF game of $F'$ with advantage  $|p_0-p_c|$. 
	
	Hint: $\B$ has to figure out how to consistently answer RO queries of $\Adv$. For example, if $\Adv$ asks $\B$ for $H(m)$, $\Adv$ may later ask $\B$ for $F'(H(m))$ (or $f'(H(m))$, depending on which game it's in). $\B$ should prepare for this eventuality, because it can't ask its wPRF challenger for a particular input. (This is called ``preparing'' or ``programming'' the random oracle.)	
	\end{ppart}
	
\begin{solution}
For this proof, let $\mathcal{B}$ be the adv. that plays against the wPRF challenger for $F'$.
It will use $\mathcal{A}$ to do this.
Note that the concept of lazy programming (from Boneh-Shoup 11.5.1) will be key in this proof.\\

\noindent
$\mathcal{B}$ will hold two associative arrays to ensure consistency between $H$ and the PRF outputs:
\begin{itemize}
    \item $\text{Map}_H$ will map msgs $m \in \M$ to inputs $x \in \X$ (basically simulating $H(m)$).
    \item $\text{Map}_F$ though, maps $x \in \X$ to outputs $y \in \Y$ (basically caching the wPRF challenger's responses).\\
\end{itemize}

\noindent
Now, $\mathcal{B}$ will run $\mathcal{A}$, and answer its queries as such:
\begin{enumerate}
    \item[RO Query $\hat{m}$]
    \begin{itemize}
      \item If $\hat{m} \in \text{Domain}(\text{Map}_H)$, return $\text{Map}_H[\hat{m}]$.
      \item Otherwise, $\mathcal{B}$ will run a query to its wPRF challenger, getting a fresh rnd pair $(x, y)$.
      Then, $\mathcal{B}$ will set $\text{Map}_H[\hat{m}] \leftarrow x$, $\text{Map}_F[x] \leftarrow y$, and then return $x$ to $\mathcal{A}$.
    \end{itemize}
    \item[PRF Query $m$]
    \begin{itemize}
      \item If $m \in \text{Domain}(\text{Map}_H)$, then let $x = \text{Map}_H[m]$, and return $\text{Map}_F[x]$ to $\mathcal{A}$.
      \item If $m \notin \text{Domain}(\text{Map}_H)$, then $\mathcal{B}$ will run a query to its wPRF challenger, getting $(x, y)$.
      $\mathcal{B}$ will then set $\text{Map}_H[m] \leftarrow x$, $\text{Map}_F[x] \leftarrow y$, and then return $y$ to $\mathcal{A}$.
    \end{itemize}
\end{enumerate}
\noindent
When $\mathcal{A}$ halts and returns $b'$, $\mathcal{B}$ will just return the same $b'$.\\

\noindent\textbf{Analysis:}
Now note, because the wPRF challenger always supplies inputs, $x$, that are uniformly rnd in $\X$, the values assigned to $H(m)$ are uniform in $\X$, which is just a true Rnd Oracle.
Consistency between $H(m)$ and the composed PRF, therefore, is guaranteed by the dual maps.

\begin{itemize}
    \item \textbf{Exp 0 (Real $F'$):} The challenger will give $(x, F'(k', x))$ for a rnd key $k'$.
    In this case, $\mathcal{A}$'s PRF queries will be answered with $F'(k', H(m))$.
    This is the same thing as the real PRF game for $\mathcal{F}(m) = F'(H(m))$.
    Thus:
    \begin{itemize}
      \item $\Pr[\mathcal{B} \to 1 \mid \text{Exp } 0] = p_0$.
    \end{itemize}
    \item \textbf{Exp 1 (Rnd Func):} The challenger now gives $(x, f(x))$ for a truly rnd function.
    In this case, $\mathcal{A}$'s PRF queries will be answered with $f(H(m))$.
    Now, this just the same thing as the behavior of Challenger $\mathcal{C}$ defined in the problem.
    Therefore:
    \begin{itemize}
      \item $\Pr[\mathcal{B} \to 1 \mid \text{Exp } 1] = p_C$\\
    \end{itemize}
\end{itemize}

\noindent By the definition of the wPRF advantage:
\[
\mathrm{wPRFadv}[\mathcal{B}, F'] = \lvert \Pr[\mathcal{B} \to 1 \mid \text{Exp } 0] - \Pr[\mathcal{B} \to 1 \mid \text{Exp } 1] \rvert = \lvert p_0 - p_C \rvert.
\]
Therefore, $\lvert p_0 - p_C \rvert \le \mathrm{wPRFadv}[\mathcal{B}, F']$, proving the OG statement of this question.
\end{solution}

\begin{ppart}
Prove that $|p_c-p_1|\le \frac{q^2}{2\cdot|\X|}$ by the difference lemma.
\end{ppart}	
	

\begin{solution}
Now this proof is basically just an extension to the previous question. I'll be making 2 more hybrid games, and then defining its ``Bad Event''.\\

\noindent\textbf{Game C ($p_C$):} The challenger will lazily sample a rnd oracle $H: \M \to \X$ and a rnd func $f': \X \to \Y$.
Upon receiving an RO query $\hat{m}$, it will return $H(\hat{m})$.
Upon getting a PRF query $m$, it returns $f'(H(m))$.\\

\noindent\textbf{Game 1 ($p_1$):} The challenger will lazily sample $H: \M \to \X$ and a rnd func $F: \M \to \Y$.
Upon getting an RO query $\hat{m}$, it returns $H(\hat{m})$.
On a PRF query $m$, it returns $F(m)$.\\

\noindent Assume that both games operate on the same underlying probability space.
In other words, the lazy sampling for $H$ will use the exact same rnd space in both games.
Now, let $W_C$ be the event that $\A$ outputs 1 in Game C, and $W_1$ be the event that $\A$ outputs 1 in Game 1.\\

\noindent Now, $\Pr[W_C] = p_C$ and $\Pr[W_1] = p_1$.\\

\noindent\textbf{$Z$:} Let $Z$ be the event that a collision in $H$ occurs on any two distinct inputs, queried by $\A$.
Since $\A$ makes AT MOST $q = q_p + q_H$ total queries, then there should be AT MOST $\binom{q}{2} < \frac{q^2}{2}$ distinct pairs of queried inputs.\\

\noindent\textbf{Analysis:} Consider the behavior of both games conditioned on $\neg Z$ occurring.
If $\neg Z$ holds, then all queried inputs $m$ should be mpped DISTINCT values in $\X$. 
\begin{itemize}
    \item In Game 1, the PRF outputs are $F(m)$, which by lazy sampling, are independent and rnd values in $\Y$.
    \item In Game C, the PRF outputs are $f'(H(m))$.
    Because $H(m)$ are all distinct and $f'$ is a rnd func, $f'(H(m))$ are also independent and rnd values in $\Y$.\\
\end{itemize}
\noindent
Thus the adversary's view in Game C and Game 1 are the exact same. This implies that $W_C \land \neg Z$ occurs IFF $W_1 \land \neg Z$ occurs.\\

\noindent\textbf{End:}
Therefore, by the difference lemma, we get:
\[
|p_C - p_1| = |\Pr[W_C] - \Pr[W_1]| \le \Pr[Z].
\]
Yet, this remains to bound $\Pr[Z]$. By the union bound over all pairs of queries, the probability that any specific pair collides is $1/|\X|$. Therefore:
\[
\Pr[Z] \le \binom{q}{2} \frac{1}{|\X|} = \frac{q(q-1)}{2|\X|} < \frac{q^2}{2|\X|}.
\]
Therefore, these inequalities yields the expected result (from the OG statement):
\[
|p_C - p_1| \le \frac{q^2}{2|\X|}.
\]

\qed
\end{solution}	
\end{problem}


\end{document}